\documentclass[UTF8,fontset=fandol,a4paper,11pt]{ctexart}
\IfFontExistsTF{Noto Serif CJK SC}{
  \setCJKmainfont{Noto Serif CJK SC}
  \setCJKsansfont{Noto Sans CJK SC}
  \setCJKmonofont{Noto Sans CJK SC}
}{
  \IfFontExistsTF{Heiti SC}{
    \setCJKmainfont{Heiti SC}
    \setCJKsansfont{Heiti SC}
    \setCJKmonofont{Heiti SC}
  }{}
}
\usepackage[margin=20mm,headheight=16pt]{geometry}
\usepackage{graphicx,xcolor,booktabs,tabularx,enumitem,fancyhdr,fvextra,etoolbox,hyperref}
\definecolor{brand}{HTML}{2080F0}
\definecolor{ink}{HTML}{25344B}
\definecolor{muted}{HTML}{607087}
\definecolor{rule}{HTML}{E3EAF4}
\hypersetup{colorlinks=true,urlcolor=brand,linkcolor=brand,pdftitle={南晓开放原子社在线报名系统配置教程},pdfauthor={南晓开放原子社}}
\graphicspath{{configuration-screenshots/}}
\setlength{\parindent}{0pt}
\setlength{\parskip}{6pt}
\setlength{\emergencystretch}{2em}
\setlist{nosep,leftmargin=1.6em,topsep=4pt}
\renewcommand{\arraystretch}{1.25}
\ctexset{section={format=\Large\bfseries\color{ink},beforeskip=0pt,afterskip=10pt}}
\pagestyle{fancy}
\fancyhf{}
\fancyhead[L]{\small\color{muted}南晓开放原子社在线报名系统}
\fancyhead[R]{\small\color{muted}配置教程}
\fancyfoot[L]{\small\color{muted}v1.0.0 · 中文 · Asia/Shanghai}
\fancyfoot[R]{\small\color{muted}\thepage}
\renewcommand{\headrulewidth}{0.3pt}
\newcommand{\code}[1]{{\small\ttfamily\detokenize{#1}}}
\newcommand{\shot}[3][0.48\textheight]{\begin{center}\fcolorbox{rule}{white}{\includegraphics[width=0.96\linewidth,height=#1,keepaspectratio]{#2.png}}\par\vspace{4pt}{\small\color{muted}#3}\end{center}}
\DefineVerbatimEnvironment{Code}{Verbatim}{fontsize=\small,breaklines=true,breakanywhere=true,samepage=true,frame=single,rulecolor=\color{rule},framesep=8pt}
\BeforeBeginEnvironment{Code}{\par\noindent\begin{minipage}{\linewidth}}
\AfterEndEnvironment{Code}{\end{minipage}\par}
\begin{document}
\thispagestyle{empty}
{\color{brand}\small 系统管理员操作手册 / 2026 年 10 月 24 日}
\vspace{9mm}

{\huge\bfseries\color{ink}南晓开放原子社\\[3mm]在线报名系统}

{\Large\color{brand}配置教程}
\vspace{5mm}

本教程从首次部署开始，说明站点与邮件设置、链接授权、报名信息导出与比赛账号、竞赛规则和准考证模板的配置方法。系统管理员维护全站配置，竞赛管理员仅管理获得授权的竞赛。

\shot[0.38\textheight]{01-home}{独立首页：点击“进入报名系统”后进入登录或竞赛列表。}

\textbf{建议配置顺序}

部署与初始化管理员 → 站点、GitHub 与邮件 → 账号资料和标签 → 链接授权 → 实名认证 → 竞赛规则与字段 → 准考证 → 测试与备份。

截图来自本地运行的实际页面。教程中的邮箱、域名、竞赛和账号均为虚构示例；截图未展示密码、插件密钥或原始邀请令牌。示例竞赛未公开发布，教程所用临时邮件设置已恢复。

\vfill
{\small\color{muted}功能范围：邮箱密码 / GitHub 登录、必填个人资料、用户标签、按时间编号审核、个人 / 团队报名、线下财务、工单、LaTeX 准考证、Excel / CSV 名单、独立比赛账号及准考证批量发放。}

\newpage
\section{首次部署与管理员初始化}
v1.0.0 按全新部署使用。已有 v0.x 实例先独立备份数据库、附件与加密密钥，再使用新的安装目录、Compose 项目名和数据卷安装；不提供旧平台账号或历史报名兼容适配。不要把旧备份直接恢复到新实例。

\textbf{一键安装。}宿主机准备 Docker、Compose 2.24.4 以上、Bash、curl 和 OpenSSL。下载根目录的统一脚本，自动下载并校验发布包；也可解压 Release 安装包后直接运行其根目录脚本：
\begin{Code}
curl -fsSL https://raw.githubusercontent.com/Albert-Li-Sz/NXOPA-RegSYS/main/install.sh -o install.sh
bash install.sh
curl --fail http://127.0.0.1:8080/api/v1/health
\end{Code}
默认目录为 \code{/opt/nxopa-regsys}，可用 \code{--dir} 指定其他目录。向导默认拉取预构建镜像，完整源码也可构建。自动生成权限为 600 的 \code{deploy/.env}，服务健康后创建首次管理员；已有配置与密钥保留，宿主机无需 Node.js 或 pnpm。内部 Caddy 提供本机 HTTP 入口，默认 8080，由自有反向代理负责公网 HTTPS。健康接口应返回 \code{status: ok}，上线前还须检查代理下的健康接口与登录。安装镜像使用完整固定版本标签，并按 Docker 服务端选择平台；新入口转换旧包的自动摘要覆盖，保留下载校验。

\textbf{更新与卸载。}不带参数显示安装、更新、卸载和状态菜单，也可在实际安装根目录分别执行：
\begin{Code}
bash install.sh update
bash install.sh uninstall
\end{Code}
更新先加密备份，再下载新版材料、拉取或构建、迁移并等待健康检查；备份失败中止操作。默认卸载保留数据，永久删除须准确输入 \code{DELETE 项目名}。配置和备份保留，备份解密需要原加密密钥。旧版 \code{install/.env} 自动复制到 \code{deploy/.env}，保留项目名、数据库密码、加密密钥与代理模式。恢复会覆盖当前数据库和附件；确认后先验证备份并建立恢复前快照。恢复失败保持业务服务停止。维护命令见 \code{bash install.sh --help}。

\textbf{设置访问地址。}向导保存下列部署项；也可编辑 \code{deploy/.env}。自行完成域名解析、证书与代理配置，参考 \code{deploy/nginx.conf.example}。首次启动使用环境地址作为初始值；登录后在“管理中心 → 全局设置”维护邮件与邀请链接使用的地址。
\begin{Code}
PUBLIC_URL=https://contest.example.org
REVERSE_PROXY_MODE=gateway
GATEWAY_PORT=8080
GATEWAY_BIND_ADDRESS=127.0.0.1
\end{Code}
外层代理将整个站点转发到 \code{http://127.0.0.1:8080}，保留路径、Host 和原始协议；内部 Caddy 分流前端与 API。默认不占用公网 80 / 443、不申请证书。旧配置升级保留原代理模式。安装入口使用 \code{PUBLIC_URL}；旧根目录 Compose 使用 \code{DEPLOY_PUBLIC_URL}。外接平台及通用开放接口已移除，GitHub 和飞书登录继续可用。

\textbf{部署环境保留：}数据库、缓存、私有附件及 \code{ENCRYPTION_KEY}。更新须保留原密钥，否则无法读取邮件密码、比赛凭据与备份。浏览器信任来源可在后台维护，环境变量保留恢复入口。

\textbf{管理中心维护：}站点、邮件、用户授权、报名导出、比赛账号、字段与准考证模板。首次登录后先配置并测试邮件。

\section{首页、登录与站点设置}
首页 \code{/} 独立展示系统介绍。未登录时点击“进入报名系统”，会跳转登录；直接访问竞赛列表或竞赛详情也会触发登录检查。登录成功后返回原本准备访问的页面，包括详情页参数。
\shot[0.31\textheight]{02-login}{登录入口：使用本站用户名或邮箱与密码。}

系统管理员登录后打开“管理中心 → 全局设置”，调整站点名称、主色与页脚，并点击“保存”。竞赛管理员没有全局设置权限。默认主色为 \texttt{\#2080F0}。
\shot[0.24\textheight]{04-site}{全局设置中的站点配置。}

\textbf{公网地址与信任来源。}公网地址用于生成邮件和邀请链接。同源请求与当前公网地址自动允许。系统管理员点击“编辑”，在“信任的请求来源”每行填写一个内外网根地址，也可用逗号分隔，保存后立即生效。例如 \code{https://reg.example.org} 与 \code{http://192.168.1.10:8080}。地址须包含协议及端口（如有），不接受通配符、账号密码、额外路径、查询参数或片段。环境变量 \texttt{TRUSTED\_\allowbreak MANAGEMENT\_\allowbreak ORIGINS} 继续可用，并在后台单独展示。保存检查版本并记录审计；TLS、域名及访问控制由反向代理负责。

\begin{enumerate}
\item 域名解析：给注册域名添加指向服务器的 A 记录；启用 IPv6 时再配置正确的 AAAA 记录。
\item 反向代理：默认由自有代理监听 HTTPS，配置证书。整个站点转发到内部 Caddy 的 \code{127.0.0.1:8080} HTTP 入口，保留 Host、Origin 和请求协议。
\item 更换域名：先解析新域名、配置代理与证书，并实际检查新站可登录，再保存新的公网地址。旧的、已经发出的邀请 URL 不会被重写，迁移期保留旧域名反向代理或引导用户重新获取链接。
\item 请求与会话：登录前亦检查写请求来源，登录后另检查 CSRF 令牌。反向代理正确传递 \texttt{X-\allowbreak Forwarded-\allowbreak Proto}。HTTPS 入口使用 Secure 会话 Cookie；公网 HTTPS 与内网 HTTP 并用时，将内网 HTTP 根地址显式加入信任来源，该入口才可使用适合 HTTP 的会话 Cookie。
\end{enumerate}
\begin{Code}
# 原生部署：在项目目录使用原数据库与加密密钥
pnpm restore:public-url https://reg.example.org
# Docker：在 API 容器内恢复数据库配置
docker compose exec api pnpm restore:public-url https://reg.example.org
\end{Code}
恢复命令保留站点其他设置、递增版本并记录审计，不需要清空数据库。


公开邮箱注册保留；用户验证邮箱并完善五项资料后才可报名，授权邀请仍需明确确认。修改登录邮箱须再认证并验证新邮箱；验证前仍使用原邮箱，完成后撤销其他会话并通知原邮箱。本站登录密码与独立比赛账号密码分别管理。

\newpage
\section{GitHub 登录与账号资料}
系统管理员打开“管理中心 → GitHub 登录”。在 GitHub 开发者设置新建 OAuth App，填写站点首页与后台显示的回调地址，再把 Client ID、Client Secret 填入本站并启用。回调路径固定为 \code{/api/v1/auth/github/callback}；自定义回调地址需要与 GitHub 配置一致。
\shot[0.48\textheight]{27-github-settings}{GitHub 登录配置：回调地址、启用状态和凭据。空白密码框表示保留已保存密钥。}

参考\href{https://github.com/nextauthjs/next-auth/blob/main/packages/core/src/providers/github.ts}{Auth.js 开源 GitHub Provider}与\href{https://docs.github.com/en/apps/oauth-apps/building-oauth-apps/authorizing-oauth-apps}{GitHub 官方 OAuth 流程}。使用一次性 state 与 PKCE；Client Secret 加密保存，GitHub 令牌不写入数据库。

\newpage
启用后登录页显示“使用 GitHub 登录”。首次授权取得 GitHub 已验证邮箱后创建本站账号；未验证邮箱不能登录。已有相同邮箱账号不会自动合并，先用本站账号登录，并在绑定前输入当前密码；没有密码时验证原邮箱验证码。只使用 GitHub 登录的账号先通过邮箱找回流程设置本站密码，再解除绑定，避免失去登录方式。
\shot[0.40\textheight]{24-github-login}{登录页同时提供本站账号与 GitHub 登录入口。}

GitHub OAuth App 的真实凭据由部署管理员配置。本版本以独立模拟 GitHub 服务测试授权、令牌交换、已验证邮箱、重复 state、绑定和同邮箱冲突，并未替管理员建立真实 OAuth App。

\newpage
\textbf{报名所需账号资料。}邮箱注册或 GitHub 登录后，在账号中心填写学号、姓名、班级、学院和性别。姓名保存在账号姓名字段，其他四项保存在个人资料；五项全部有效才允许创建或提交报名。旧账号缺项时补充后可继续报名，原报名与历史快照保留。
\shot[0.54\textheight]{26-profile}{账号中心的五项必填资料。填写完成并保存后，返回原报名页面。}

注册完成后直接引导补充资料；已登录用户从竞赛点击报名时也会检查。团队中已绑定本站账号的参赛成员分别完善资料；指导老师资料仍按竞赛规则填写。

\newpage
\section{邮件服务：保存、检测与实际收件}
打开“管理中心 → 邮件设置”。只有系统管理员可以查看、修改或测试邮件配置。未在后台保存时，使用部署环境中的初始 SMTP 配置；后台保存后优先使用数据库配置，无需重启 API 或后台任务服务。
\shot[0.47\textheight]{03-mail}{本地邮件测试示例：保存配置，检测连接，再向当前管理员发送测试邮件。}

\begin{tabularx}{\linewidth}{@{}lX@{}}\toprule
配置项 & 填写方法\\\midrule
服务器地址 / 端口 & 填服务商的 SMTP 主机名，不含协议与路径；端口按服务商要求填写。\\
连接加密 & 465 通常用“加密连接”；587 通常用“连接后升级加密”。按服务商文档确认。\\
服务器认证 & 通常启用。填写服务商用户名及密码或邮箱授权码。\\
发件邮箱 / 名称 & 使用服务商允许的发件地址；名称建议填写系统名称。\\\bottomrule
\end{tabularx}

先保存，再“检测连接”，最后“发送测试邮件”。测试邮件只发给当前管理员。任务成功后还应检查实际收件箱或垃圾邮件；连接检测通过不能单独证明邮件可送达。加密方式对应关系见\href{https://nodemailer.com/smtp}{官方邮件配置说明}。

密码加密保存且不返回浏览器。编辑时留空保留原密码；更换主机、端口、加密方式或认证用户名时重新填写密码。关闭服务器认证并保存，会清除认证密码。关闭邮件服务后不能发送验证和找回邮件，失败任务会保留原因供管理员处理。

\newpage
\section{通过链接授予账号身份}
打开“管理中心 → 用户权限”，点击“生成邀请链接”。选择邀请类型、已有账号或受邀邮箱；竞赛管理员还需选择授权竞赛。默认有效期为 7 天，可设置 1 至 30 天。
\shot[0.55\textheight]{05-role-invitation}{生成竞赛管理员授权链接的配置表单。}

\begin{enumerate}
\item 点击“生成链接”并复制 URL，交给指定接收者。链接只在生成时展示。
\item 接收者通过该 URL 登录或注册指定邮箱，完成邮箱验证。
\item 接收者在邀请页面查看身份和竞赛范围，点击“确认接受”后授权才生效。
\end{enumerate}
仅生成或打开链接不会增加权限。账号邀请不附加管理员权限；系统管理员、竞赛管理员、指导老师身份分别邀请。撤销未确认邀请立即使链接失效；撤销已有角色后，旧链接不能重新恢复权限。

\newpage
\textbf{用户标签与筛选。}在“用户权限”列表点击“编辑标签”，输入或选择标签并保存，例如“集训队”“2026级”。每个账号最多 20 个标签，单个标签最多 40 字；重复值自动合并。标签只用于管理筛选，不授予权限。
\shot[0.48\textheight]{25-user-tags}{管理员用户列表展示标签，并可按标签筛选。}

用户列表可按标签筛选；竞赛管理员可在报名列表按成员标签筛选。列表、统计与 CSV 导出采用相同筛选条件。更新标签检查版本，冲突时刷新后重新保存。

\newpage
\section{授权确认与范围检查}
\shot[0.46\textheight]{06-role-confirmation}{接收者确认前查看身份、受邀邮箱、有效期和指定竞赛。}

确认后刷新管理中心，竞赛管理员应只看到获得授权的竞赛，并可处理对应报名、财务、工单和准考证。竞赛管理员不能修改全站 SMTP 设置或评测实例的接入凭据。

指导老师身份也通过邀请确认。老师与某份报名的关联还需通过链接确认，关联解除后立即失去该报名的访问和代办权限。队伍成员邀请保持独立确认流程；竞赛选择“队长代填”时按照竞赛规则录入人员信息。

\textbf{排查常见问题}
\begin{itemize}
\item 邮箱不匹配：使用邀请指定邮箱登录，不能由其他账号接受。
\item 邮箱未验证：到账号中心重新发送验证邮件，再返回邀请页面。
\item 链接过期或撤销：由系统管理员生成新的邀请。
\item 权限范围不符：检查是否选中了正确竞赛，以及身份是否已经被撤销。
\end{itemize}

\newpage
\section{导出 Excel 与 CSV 报名信息}
在活动的“报名管理”中点击“导出名单”，选择 Excel（.xlsx）或 CSV。默认使用 Excel，导出范围覆盖当前筛选条件下的全部报名，每名选手一行。团队资料重复到各成员行，包含成员、老师、实名单位、动态字段、审核、资格、考点、考场与座位；财务列要求财务权限，金额只在每份报名的首行填写。

Excel 冻结表头并提供筛选，学号、手机号和账号按文本保存，保留前导零和长编号。CSV 使用带 BOM 的 UTF-8，并转义公式字符。每次最多导出 10000 名成员、500 列和 25 MiB 文件；请用筛选条件缩小超大名单。

\section{配置独立比赛账号}
活动设置中的“生成比赛账号”默认关闭。开启后填写账号前缀，例如 \code{nxcpc_}，账号按报名编号生成，如 \code{nxcpc_0001}。个人报名一组，团队全队共享一组，密码随机生成并加密保存，与本站登录密码分开。

报名取得参赛资格后自动生成；中途开启功能会补齐已有合格报名。重复生成、导出或关闭再开启都保留原密码。首个账号生成后前缀锁定。资格失效或关闭功能后停止输出密码，恢复资格后复用原账号。

“准考证显示账号密码”默认开启，模板变量为 \code{account.username} 和 \code{account.password}。导出弹窗另有“包含比赛账号密码”，默认不勾选。修改账号显示配置后，旧准考证失效，需要重新生成、检查和发放。

\section{实名单位与批量准考证}
学校／单位是实名资料必填项。个人报名引用本人、团队报名引用队长审核通过的单位，只读显示；不再单独填写单位简称、ICPC 编号或国家／地区代码。可以保存报名草稿，但提交及补正重提时，本人／队长必须已经通过实名审核且单位完整。团队允许跨单位成员，代表单位在提交时写入报名快照。

报名管理支持跨页勾选，也可显式选择当前筛选条件下的全部报名。改变筛选条件会清空勾选，翻页保留勾选。操作前预览总数、可处理数与跳过原因，并固定目标报名 ID 和版本。

批量生成完成后，管理员先检查 PDF，再单独执行“批量发放准考证”。发放后开放站内下载并触发通知，重复发放不重复通知。后台显示成功、失败、跳过及进度，资料变化或资格失效的记录不会继续处理。每批最多 10000 份；ZIP 每包最多 1000 份 PDF、源文件合计最多 100 MiB，超出自动分包，下载入口 24 小时有效并重新校验全部文档权限。

\newpage
\section{竞赛报名规则与线下财务}
打开“活动及竞赛管理中心”，创建或查看活动。已保存配置先点击“编辑”才可修改；设置名称、赛季、时间、地点、公告和发布状态，再配置报名规则。草稿不显示在公开列表。
\shot[0.30\textheight]{10-contest-policy}{团队报名示例：邀请成员、人数限制、老师关联与人工审核。}

个人报名固定一人。团队报名设置最少 / 最多人数，以及成员邀请或队长代填。老师可设置为必需、可选或不启用；必需老师的报名需要有效确认关联。
\shot[0.28\textheight]{11-contest-finance}{线下财务示例：按报名单计费，先审核再缴费。}

可设置免费、按报名单或按人数计费；费用输入单位为元，支持两位小数，例如 150.00 元。填写线下收款说明，选择审核与缴费顺序，并设置退款与开票。非正式报名同样需要审核。核账只能确认凭证申报金额，金额不符须退回；应记录交易参考号，申请人和上传者不能核账自己的凭证。部分到账不视为足额缴费，超应收退款须由不同管理员批准和完成。

\newpage
\section{配置考点与考场}
在竞赛配置页点击“编辑”，在“考点选择”中新增考点，填写名称和地址，设置是否开放选择后保存。每个竞赛最多 100 个考点，名称不可重复。未配置考点时，报名沿用原流程。

配置考点后，选手提交报名须选择一个开放考点，团队成员共用同一考点。草稿可以暂不选择，但提交前必须补充；停用考点不再接受新选择，已提交报名保留原选择。选手更换考点后须重新提交审核，旧座位会清除，旧准考证会失效。

已被报名或考场使用的考点不能直接删除，可先停用，或调整关联数据后再删除。已有报名的竞赛后期新增考点时，须先补充报名考点，再排考或生成准考证。

启用“考场与座位”后，为每个考场选择所属考点。自动、手动及 CSV 排座均限制在报名所选考点内；团队按现有排考模式安排同考点座位。考场已有座位时，先清空座位才能更换所属考点。准考证地点合并考点名称、地址和考场地址；修改考点名称或地址后，应重新生成并发放准考证。

考场另有可选的“考场编号”，支持字母、数字、下划线和连字符，最长 20 字符，去除首尾空格后统一大写，同一活动内唯一。报名导出分别包含考场与座位。修改编号、地点和排考后，相关准考证需要重新生成。

管理员在报名管理中可查看和筛选考点，导出的名单包含考点名称与地址。报名详情也显示所选考点及地址，便于选手核对。

\newpage
\section{配置报名字段}
\textbf{按时间审核。}每个活动从 1 开始按首次成功提交的先后分配报名编号。草稿没有编号；并发提交不重号；补正与重新提交保留原编号与首次提交时间。取消后永久删除报名，财务独立保留，原编号不重复使用。名单 CSV 为每位有效成员一行，每个成员、队伍、报名和老师字段各一列；团队金额仅填写在该报名第一行，单位为元。
在竞赛“报名字段”页增加字段，或从全局字段导入。设置中文名称、稳定键名、作用域、类型、必填与敏感标记，以及选项或条件。提交报名会保存人员、表单与费用快照；关键资料修改需要重新审核。
\shot[0.38\textheight]{12-fields}{可视化字段列表：参赛者、队伍、老师或报名单的字段分别配置。}

\begin{tabularx}{\linewidth}{@{}lX@{}}\toprule
作用域 & 使用场景\\\midrule
参赛者 & 学校、学号、专业、电话等个人资料。\\
队伍 & 队伍标识、队伍附加资料。\\
指导老师 & 老师关联后补充填写的竞赛信息。\\
报名单 & 整份报名需要提交的其他信息。\\\bottomrule
\end{tabularx}

键名用于接口和模板变量，保存后保持稳定。类型使用短文本、长文本、数字、日期、单选、多选、开关、文件或照片。选项字段需要提供选项内容；条件字段只在条件满足时显示。敏感字段在接口查询中需要明确授权。

报名字段按参赛者、队伍、报名和老师划分，支持自定义字段与历史版本。代表单位直接引用已通过审核的实名单位，平台基础模板已移除。

\newpage
\section{发票申请与审核}
参赛者在报名“财务”页提交完整开票表单，先查看可开票余额并确认信息。抬头分个人、企业、非企业单位；默认仅普通发票，管理员在竞赛费用配置中开启专用发票并选择额外必填资料。
\shot[0.45\textheight]{21-invoice-application}{完整开票申请表：票种、抬头、税号、金额、联系方式、银行资料和备注。}

企业税号必填；个人和非企业单位的普通发票允许留空。专票必须为单位抬头并填写税号；地址、银行等是否必填由开票方配置。不要把所有票种都强制成传统纸质专票字段，参考\href{https://shanghai.chinatax.gov.cn/tax/zcfw/rdwd/202402/t470676.html}{税务机关购买方信息说明}。
申请金额不能超过确认净收款减去已申请／已开票金额；退款前须处理关联发票。提交保存资料快照，财务管理员及报名参与者可以查看完整信息，普通查看或报名审核授权仅能查看费用摘要。拒绝时填写原因，实际开票或退款时绑定凭证；历史财务附件继续保留。历史申请缺少新字段时显示为历史记录，不补造资料。

\newpage
\section{分页、Markdown 与页面使用}
附件上传期间先等待完成，再保存、取消或退出登录。会话过期时可使用原账号恢复登录，当前草稿会保留。

管理中心的用户、授权邀请、审计与任务，报名名单、工单以及财务历史均支持 20／50／100 条每页。显示总数、页码与跳页；页大小刷新后保留，筛选变化回到第一页。搜索、状态与时间条件同时作用于列表和总数；统计和名单导出始终覆盖完整筛选结果。批量准考证操作使用跨页勾选或显式选择的完整筛选结果，并在执行前确认数量。
\shot[0.29\textheight]{22-pagination}{中文审计操作与分页：选择每页条数，查看关联对象和详情。}

竞赛介绍可选择纯文本或 Markdown。旧介绍保持纯文本；编辑器提供标题、列表、表格、引用、代码块、链接和图片地址工具与预览。渲染关闭原始 HTML 并净化内容；图片使用可公开访问的 HTTPS 地址。手机号、税号等私人信息不要写入公开介绍。
\shot[0.31\textheight]{23-markdown}{竞赛 Markdown 编辑与预览，支持表格、引用及代码块。}

\newpage
\section{配置准考证模板与发放}

准考证遵循个人报名一人一张、团队报名全队一张。系统管理员在“全局设置”维护全局模板；竞赛管理员在竞赛“准考证模板”页保存竞赛版本、选择字段并预览。
\shot[0.60\textheight]{13-admission-template}{模板源码、保存与预览按钮；右侧字段面板可插入变量。}

竞赛模板优先于全局模板。管理员在报名管理中选取已获资格的报名，配置考点、考场、座位，批量生成并确认发放。用户只可下载已发放且版本有效的文件；资料、资格或模板变化会使旧文件失效。

\newpage
\section{模板语法、任务与维护}
模板仅使用受限变量、条件与循环。字段数据统一转义，不能直接插入任意原始代码。下面展示团队成员循环；个人字段可使用 \code{participant.name} 与 \code{participant.answers.school}。
\begin{Code}
\documentclass[UTF8,fontset=fandol]{ctexart}
\begin{document}
{{contest.title}} 准考证
编号：{{admission.number}}
考场：{{admission.room}}  座位：{{admission.seat}}
{{#each team.members}}
{{this.name}}：{{this.answers.school}}\\
{{/each}}
\end{document}
\end{Code}
预览无误后保存新版本。生成任务失败时查看编译原因，修复后重新生成；失败文件不会向参赛者发放。完整变量、照片字段和座位导入格式见 \code{docs/ADMISSIONS.md}。
\shot[0.12\textheight]{14-jobs}{任务页：检查后台任务服务、待派发事件和失败任务。}

在“任务”页查看状态、错误和尝试次数，失败任务可重试。邮件测试状态表示服务器已接收邮件，实际送达仍需检查收件箱。批量准考证任务显示逐项进度和结果，权限或资料变化后需要重新核对。网站内部接口说明见 \code{docs/API.md}。

\begin{Code}
bash scripts/backup.sh /srv/contest-backups
bash scripts/restore.sh /srv/contest-backups/contest-center-YYYYMMDDTHHMMSSZ.ccbak
\end{Code}
升级前备份数据库、私有附件与环境文件；备份包含后台邮件设置。异机恢复使用原加密密钥，并先在隔离环境核验。数据库连接等基础设施变更仍需更新部署环境并重启服务；公网地址与邮件配置在后台保存后生效。环境配置的可信管理入口保留，以便恢复错误公网地址。
\end{document}
